\originalchapter{辐射自由度、中微子温度与热史的分段}
本章补齐8.286第15--18讲的量子气体和中微子内容。热遗迹章已推导黑体能量、数密度、熵和冻结；这里解释不同统计权重怎样从积分产生，以及脱耦后为什么不能把所有粒子放进同一个温度。课程2013年的中微子参数和历史CMB谱图只作年代明确的背景，原第三方图片不收录。

\originalsection{同样多的自旋态为什么贡献不同的能量}
先取$\hbar=c=\kb=1$。具有$g$个内部态的自由粒子，相空间计数为$g\dd^3p/(2\pi)^3$。若化学势可忽略，平衡分布是$f_B=(\ee^{E/T}-1)^{-1}$或$f_F=(\ee^{E/T}+1)^{-1}$。对无质量粒子，$E=p$，故
\[
 n=\frac{gT^3}{2\pi^2}\int_0^\infty\frac{x^2\dd x}{\ee^x\mp1},\qquad
 \epsilon=\frac{gT^4}{2\pi^2}\int_0^\infty\frac{x^3\dd x}{\ee^x\mp1}.
\]
用$1/(\ee^x-1)=\sum_{j\geq1}\ee^{-jx}$，逐项积分$\int x^{s-1}\ee^{-jx}\dd x=\Gamma(s)j^{-s}$；费米展开则有交替符号。于是
\[
 \int_0^\infty\frac{x^{s-1}\dd x}{\ee^x+1}
 =(1-2^{1-s})\Gamma(s)\zeta(s).
\]
能量用$s=4$，权重$7/8$；数目用$s=3$，权重$3/4$。它们来自不同积分，不是可以随意互换的“费米修正”。压强从各向同性动量流$\int p_iv_i f\dd^3p$给$\epsilon/3$，熵$ s=(\epsilon+p)/T$因而与能量有同样的$7/8$权重。

原8.286把能量／熵有效数记作$g$，把数目有效数记作$g^*$；本卷采用$g_*$表示能量、$g_{*s}$表示熵，另记$g_{*n}$表示数目。不同温度物种贡献分别带$(T_i/T)^4$、$(T_i/T)^3$、$(T_i/T)^3$。符号转换后关系为
\[
 \epsilon=\frac{\pi^2}{30}g_*T^4,\quad
 s=\frac{2\pi^2}{45}g_{*s}T^3,\quad
 n=\frac{\zeta(3)}{\pi^2}g_{*n}T^3.
\]
质量不小于温度时需要保留$E=\sqrt{p^2+m^2}$积分，不能以阶跃的自由度表作精确热史。

\originalsection{MeV时期的自由度为什么是10.75}
光子有两个横向极化。电子和正电子各有两个自旋态，在$T\gg m_e$时共四个费米自由度。标准热背景中每个活跃中微子味有一个中微子螺旋态和一个反中微子螺旋态，三个味合计六个，所以
\[
 g_*=2+\frac78\,4+\frac78\,6=\frac{43}{4}=10.75,
 \qquad g_{*n}=2+\frac34\,4+\frac34\,6=9.5.
\]
这里假定这些物种都在近共同温度，且未额外热化其他态。中微子有质量，不意味着Dirac模型的其他螺旋态必定在早期热浴中以同样丰度产生；相互作用率仍须检查。Majorana与Dirac的粒子解释不同，但通常的活跃热中微子计数在上述近似下相同。

四费米弱作用在$T\ll m_W$时给数量级$\sigma\sim G_F^2T^2$，背景数密度$\sim T^3$，故$\Gamma\sim G_F^2T^5$。膨胀率$H\sim T^2/M_P$；随降温$\Gamma/H\propto T^3$下降，交叉在MeV量级。精确脱耦不是一个瞬间，三个味也非完全同步；这里的突变近似用于解释主效应。

\originalsection{电子湮灭怎样使光子比中微子热}
中微子脱耦后，其动量分布随$a^{-1}$平移，温度参数$T_\nu\propto a^{-1}$。电子与正电子仍和光子耦合，湮灭时将熵主要传给这一部门。湮灭前该部门的熵自由度是$2+7\times4/8=11/2$，湮灭后为2。若两端充分相对论／完全湮灭且无额外熵交换，部门熵守恒给
\[
 \frac{11}{2}(a_iT_i)^3=2(a_fT_{\gamma,f})^3,
 \qquad a_fT_{\nu,f}=a_iT_i,
 \qquad \frac{T_\nu}{T_\gamma}=\left(\frac4{11}\right)^{1/3}.
\]
因而$T_\gamma=2.725$ K时$T_\nu\simeq1.945$ K。这个“温度”在中微子变成非相对论性后是描述遗留动量分布的参数，不能再把它说成与同温物质反复碰撞形成的平衡温度。

每个味含粒子和反粒子的数密度相对光子为
\[
 \frac{n_{\nu_i}+n_{\bar\nu_i}}{n_\gamma}
 =\frac{(3/4)\,2}{2}\left(\frac4{11}\right)=\frac3{11}.
\]
光子$n_\gamma\simeq411\ \mathrm{cm^{-3}}$，所以每味约112$\ \mathrm{cm^{-3}}$。若仍相对论性，三味能量使
\[
 \frac{\epsilon_r}{\epsilon_\gamma}
 =1+\frac78\left(\frac4{11}\right)^{4/3}3\simeq1.681.
\]
实际有限时长的脱耦及有限温度修正通常用$N_{\rm eff}$写为$1+(7/8)(4/11)^{4/3}N_{\rm eff}$，而不是把非整数$N_{\rm eff}$当作已经发现了非整数个粒子种类。本卷不在突变模型中冒充已算出高精度修正。

\originalsection{中微子质量与背景动力学的联系}
当$p\ll m_\nu$时，$\epsilon_\nu\simeq\sum_i m_i n_i$，以标准遗留数密度除以今日临界密度可得近似
\[
 \Omega_\nu h^2\simeq\frac{\sum_i m_ic^2}{94\ \mathrm{eV}}.
\]
系数取决于标准温度和脱耦约定。这不是质量测量；它把某个假定质量和热史转成宇宙密度。振荡测得质量平方差，不能独自确定绝对质量和最轻本征态。第18讲所列约$7.5\times10^{-5}$与$2.3\times10^{-3}\ \mathrm{eV^2}$为2013年例值，不作当前最佳拟合声明。

非相对论背景中微子不应同时算在$a^{-4}$的辐射密度和$a^{-3}$的物质密度里。精确转换期应计算
\[
 \epsilon_\nu(a)=\frac{g}{2\pi^2a^4}\int_0^\infty
 q^2\sqrt{q^2+a^2m^2}\,f(q)\dd q,
\]
其中$q=ap$为共动动量。高温与低温极限分别给两种尺度律，过渡连续。自由流还能抑制小尺度物质增长，密度记账与扰动效应是两个相关却不同的计算。

\originalsection{时间、温度与可见的历史证据}
在$g_*=10.75$近不变的辐射时期，$t\simeq0.738(T/\mathrm{MeV})^{-2}$ s，即$T\simeq0.859/\sqrt{t/\mathrm s}$ MeV，与第18讲的0.860一致。电子湮灭后$g_*$与光子温度关系改变，需要另算系数；复合附近又已不是纯辐射主导。把一个温度—时间式从MeV一直外推到今天会混淆这些阶段。

第18--19讲回顾1975年的早期谱数据、1987年火箭结果与1990年COBE/FIRAS黑体谱。其物理意义是宽频谱与同一黑体温度相容，不是所有早期测量都同等可靠。本卷的Planck曲线由公式生成，不复制原测量图，也不编造原误差棒。黑体\emph{平均谱}支持热起源，角方向的微小\emph{涨落谱}检验几何和结构种子；后者须联系CMB章的$C_\ell$，不能用前者的温度拟合取代。
